% Transcrito de Documento_Visao_MotivaFit_Completo_e_Detalhado.pdf
\clearpage
\section*{REFERÊNCIAS}
\addcontentsline{toc}{section}{REFERÊNCIAS}

[1] MOORE, Geoffrey A. Crossing the chasm: marketing and selling high-tech products to mainstream customers. New York: Harper Business, 1991.

[2] MARCZEWSKI, Andrzej. User Types. In: MARCZEWSKI, Andrzej. Even ninja monkeys like to play: gamification, game thinking and motivational design. 1. ed. [S.l.]: CreateSpace Independent Publishing Platform, 2015. p. 65--80.

[3] TONDELLO, Gustavo F. et al. The gamification user types Hexad scale. In: ANNUAL SYMPOSIUM ON COMPUTER-HUMAN INTERACTION IN PLAY, 2016, Austin. Proceedings [...]. New York: ACM (Association for Computing Machinery), 2016. p. 229-243. DOI (Digital Object Identifier, identificador digital de objeto): 10.1145/2967934.2968082.

[4] PRESSMAN, Roger S.; MAXIM, Bruce R. Engenharia de software: uma abordagem profissional. 8. ed. Porto Alegre: AMGH, 2016.

[5] SCHWABER, Ken; SUTHERLAND, Jeff. O Guia do Scrum: O Guia Definitivo para o Scrum: As Regras do Jogo. 2020.

[6] BASILI, Victor R.; CALDIERA, Gianluigi; ROMBACH, H. Dieter. The goal question metric approach. In: MARCINIAK, John J. (ed.). Encyclopedia of software engineering. New York: John Wiley \& Sons, 1994. v. 1, p. 528-532.

[7] KRUCHTEN, Philippe; NORD, Robert L.; OZKAYA, Ipek. Managing technical debt: reducing friction in software development. Boston: Addison-Wesley, 2019.

[8] BRASIL. Ministério do Planejamento, Orçamento e Gestão. Guia de contagem de pontos de função do MP (Ministério do Planejamento). Versão 1.0. Brasília, 2015. Disponível em: \url{https://bibliotecadigital.gestao.gov.br/handle/777/235}. Acesso em: 30 set. 2026.

[9] HEVY. Gym progress. Disponível em: \url{https://www.hevyapp.com/features/gym-progress/}. Acesso em: 30 set. 2026.

[10] STRONG. Workout tracker and gym log. Disponível em: \url{https://www.strong.app/}. Acesso em: 30 set. 2026.

[11] JEFIT. Workout logging app. Disponível em: \url{https://www.jefit.com/use-case/workout-logging-app}. Acesso em: 30 set. 2026.

[12] AMERICAN SOCIETY FOR QUALITY (ASQ). Fishbone diagram. [S.l.] (sem local), [s.d.] (sem data). Disponível em: \url{https://asq.org/quality-resources/fishbone}. Acesso em: 30 set. 2026.
